\chapter{When a random invitation still leaves selective recovery}
\label{rel73:partial-recovery}

A bank's plan to recover missing customer histories can fail at the follow-up stage. A bank may instead invite a random sample
of customers whose records are incomplete and receive usable evidence from
only some of them. Customers with a resolved dispute may reply more readily
than those still in difficulty; a collections call can itself alter a later
repayment. Randomizing the invitation does not make the completed responses
representative, nor does it make a later outcome equal to the historical
outcome the bank intended to recover. Both distinctions change the conclusion
that the bank can draw from the follow-up expense.

Consider a simpler decision with a genuinely bounded endpoint. A servicing
offer is randomized before a fixed repayment date. The endpoint is whether
the specified payment was made by that date, and the target is the average
effect of the offer in the experimental customer population. The bank later
tries to recover the historical payment status from incomplete records. This
example concerns a binary historical outcome and a single treatment. It does
not impose artificial finite bounds on unbounded Gaussian servicing rewards
or claim to recover a longitudinal policy value from a terminal indicator.

\section{The information supplied by an invitation}

Let $A\in\{0,1\}$ denote randomized treatment and $Y(a)\in\{0,1\}$ the
historical outcome under assignment $a$. Consistency and no interference give
$Y=Y(A)$; random assignment identifies the population contrast
$\tau=E\{Y(1)-Y(0)\}$ from the two full arm means if those outcomes were
observed. Let $R$ indicate that routine records contain $Y$. For customers
with $R=0$, let $B$ indicate an invitation and let $S$ indicate whether the
specified follow-up would recover an outcome if that customer were invited.
An uninvited customer's $S$ is not observed. The design requires
\[
 B\mathbin{\perp}(Y,S)\mid A=a,R=0,\qquad
 \Pr(B=1\mid A=a,R=0)=q_a>0.
\]
The invitation probabilities are known. The outcome among completed
follow-ups must be the same historical $Y$, not a new measurement of a
changed endpoint. Neither routine response nor follow-up completion is
assumed independent of $Y$.

For each arm define the routine success mass, recovered success mass and
remaining mass by
\[
 v_a=E(RY\mid A=a),\qquad
 z_a=E\{(1-R)SY\mid A=a\},\qquad
 m_a=E\{(1-R)(1-S)\mid A=a\}.
\]
Random invitation identifies $z_a$ and $m_a$ through the observed conditional
distribution of follow-up response and outcomes among invitees. It does not
identify the outcome mean in the remaining mass. Binary support consequently
gives arm bounds $\ell_a=v_a+z_a$ and $u_a=\ell_a+m_a$, and effect bounds
\begin{equation}
 I=[\tau_L,\tau_U]=[\ell_1-u_0,\ u_1-\ell_0],\qquad
 \tau_U-\tau_L=m_1+m_0.
 \label{rel73:identified-set}
\end{equation}
This is the binary specialization of the double-sampling bounds developed
by \citet{rel73coppock2017}. Their outcome-stability assumption permits
arbitrary differences between initial respondents and nonrespondents. It
restricts the meaning of the recovered outcome, not the selectivity of the
initial response.

The bounds are attainable within the declared model. Assign outcome zero
to every unresolved treated customer and one to every unresolved control
customer to obtain the lower endpoint; reverse those assignments for the
upper endpoint. These assignments leave the routine observations,
invitation law, completion probabilities and observed follow-up outcomes
unchanged. The joint coupling of the two potential outcomes is otherwise
unrestricted, so the arm completions can coexist in a population. This
establishes sharpness relative to these observed data and assumptions; extra
restrictions on outcomes or selection would define a different model.

Holding the follow-up procedure fixed, changing a positive invitation
probability does not change $m_a$ or the population identified set. It
changes how precisely the available follow-up law can be estimated. A more
effective recovery procedure can reduce $m_a$ and hence identification width.
At $q_a=0$, the follow-up law is not identified in that arm, so the conclusion
cannot be extrapolated to no invitations. If increasing invitations changes
staff attention, response behavior or outcomes, it also changes the procedure;
the fixed-law comparison no longer applies.

\section{Estimating a set and covering it are separate operations}

For independent customer observations across and within the two arms,
and invitations drawn independently across customers, known Bernoulli
invitation permits
the following observable scores, with missing products evaluated only for
records whose outcomes are observed:
\begin{equation}
 L=RY+\frac{(1-R)BSY}{q_a},\qquad
 U=RY+(1-R)\left\{1-\frac{BS(1-Y)}{q_a}\right\}.
 \label{rel73:endpoint-scores}
\end{equation}
Conditional expectation over $B$ gives $E(L\mid A=a)=\ell_a$ and
$E(U\mid A=a)=u_a$. Thus
$\widehat\tau_L=\overline L_1-\overline U_0$ and
$\widehat\tau_U=\overline U_1-\overline L_0$ are unbiased endpoint
estimators. These are authored inverse-invitation estimators under Bernoulli
sampling. They are not a native execution of the published package's
sample-ratio estimators or its fixed-follow-up-size variance calculation.

Inverse weighting can give an individual score outside the outcome support.
Here $L\in[0,1/q_a]$ and $U\in[1-1/q_a,1]$, so each range has length
$1/q_a$. In a finite sample the two estimated endpoints can cross. Sorting
them or concealing such a sample would change the declared procedure. The
raw estimates therefore remain in the results, along with any empty
confidence set.

A conservative finite-sample guarantee follows from the bounded-sum inequality of \citet[Theorem 2]{closure30hoeffding}, applied to these score ranges. Let $n_a$ be the number of independent customers in arm $a$, and put
\begin{equation}
 D=\frac{1}{n_1q_1^2}+\frac{1}{n_0q_0^2},\qquad
 t_\alpha=\sqrt{\frac{D}{2}\log\frac{2}{\alpha}},\qquad
 C_\alpha=[\widehat\tau_L-t_\alpha,
            \widehat\tau_U+t_\alpha]\cap[-1,1].
 \label{rel73:confidence-set}
\end{equation}
Under the stated independence, design and stability assumptions,
$\Pr(I\subseteq C_\alpha)\geq1-\alpha$. The guarantee covers the entire
population identified interval, and consequently the true effect within it.
It does not say that the reported set is narrow enough to support an action.

For completeness, the argument needs no estimated variance. If a random
variable has range length $c$, the second derivative of its centered
log-moment-generating function is its variance under exponential tilting.
That variance is at most $c^2/4$: the squared distance from the range
midpoint is bounded by $c^2/4$, and the variance is no larger. Integrating
twice from zero gives a log-moment bound of $s^2c^2/8$. Independence then
bounds either centered endpoint estimate's log moment by $s^2D/8$.
Exponential Markov inequality, minimized at $s=4t/D$, bounds either adverse
tail by $\exp(-2t^2/D)$. Applying this to overestimation of the lower endpoint
and underestimation of the upper endpoint, and summing the two failure
probabilities, proves~\eqref{rel73:confidence-set}. Dependence between the
two endpoint estimates is allowed; independence of the customers used in
the product argument is not optional. Intersection with $[-1,1]$ preserves
coverage because the population identified interval is inside that support.

The comparison also reports normal endpoint bands using sample variances of
the scores and a 1.96 multiplier on each outward standard error. Their
behavior is empirical here; the preceding finite-sample guarantee does not
transfer to them. The original article's Imbens--Manski construction instead
targets asymptotic coverage of the true effect, which is a different
criterion from simultaneous containment of both endpoints. A software label
such as \emph{confidence interval around the region} does not erase this
distinction. \citep{rel73coppock2017,rel73attrition}

\section{Selective completion can reverse a point conclusion}

The retained generated experiment has historical arm means .45 and .55,
so the true effect is .10. Routine observation probabilities, indexed by
outcome zero and one, are $(.8,.4)$ in control and $(.4,.8)$ in treatment.
Both arms therefore have routine response probability .62, but their observed
success means differ sharply. Equal response rates do not imply equal
selection. The routine-only complete-case contrast is approximately .4194,
while routine-only worst-case bounds are $[-.12,.64]$.

Three stable recovery procedures vary successful completion among invited
nonresponders. The standard procedure has probabilities $(.1,.9)$ in control
and $(.9,.2)$ in treatment; the intensive procedure has $(.8,.98)$ and
$(.98,.8)$ respectively; complete recovery has probability one in both
arms. These probabilities define the generated truth. The estimator is not
given the outcome-dependent selection law or the true arm means.

\begin{table}[htbp]
\centering\small
\begin{tabular}{lrrr}
\toprule
Recovery procedure & Population effect set & Width & Naive limit\\
\midrule
Routine records only & $[-.1200,.6400]$ & .7600 & .4194\\
Standard & $[-.0870,.1540]$ & .2410 & $-.0720$\\
Intensive & $[.0560,.1108]$ & .0548 & .0697\\
Complete & $[.1000,.1000]$ & .0000 & .1000\\
\bottomrule
\end{tabular}
\caption{Exact generated population quantities. The routine-only naive
contrast uses complete cases. The other naive limits replace every routine
nonrespondent's outcome by the mean among completed follow-ups in that arm.
All stable sets contain the true effect .10.}
\label{rel73:population}
\end{table}

Under standard recovery, successful control follow-ups overwhelmingly reveal
historical successes, whereas successful treated follow-ups disproportionately
reveal failures. Treating those replies as representative produces the
negative limiting contrast in Table~\ref{rel73:population}. A random
invitation has not repaired selective completion. The valid set retains
both signs; the intensive procedure supplies enough extra outcome information
to exclude zero at the population level.

Each procedure is evaluated in 500 independent datasets with 2,000 customers
per arm. A fourth, deliberately unstable procedure is discussed below.
Within a dataset, shared invitation uniforms compare probabilities
.10, .25, .50 and 1.00. There are 2,000 independent generated datasets and
8,000 procedure--invitation-rate cells, not 8,000 independent datasets.
All cells remain in the denominator. Reaggregation checks every retained
endpoint, interval indicator, resource total and summary; direct observed-cell
sums and squares reconstruct endpoint estimates and standard errors from
the first, middle and last retained raw datasets in each procedure.

\begin{table}[htbp]
\centering\small
\begin{tabular}{llrrrr}
\toprule
Procedure & Invitation & \multicolumn{2}{c}{Whole-set coverage}
 & \multicolumn{2}{c}{Mean set width}\\
 & probability & Bounded & Normal & Bounded & Normal\\
\midrule
Standard & .10 & 100.0\% & 97.6\% & 1.100 & .376\\
Standard & 1.00 & 99.2\% & 96.0\% & .327 & .302\\
Intensive & .10 & 100.0\% & 96.6\% & .915 & .220\\
Intensive & 1.00 & 99.4\% & 94.8\% & .141 & .117\\
Complete & .10 & 100.0\% & 98.0\% & .860 & .172\\
Complete & 1.00 & 99.2\% & 95.4\% & .086 & .062\\
\bottomrule
\end{tabular}
\caption{Selected endpoints of the retained invitation comparison, with
500 replications per row and nominal 95\% confidence. Bounded refers to
equation~\eqref{rel73:confidence-set}; Normal is the separately assessed
normal band. Coverage requires containment of the entire population set.
An empty confidence set has width zero and fails coverage.}
\label{rel73:coverage}
\end{table}

The bounded construction buys a finite-sample statement with substantial
conservatism. Under intensive recovery at invitation probability .10 it
declares no positive effects, even though the population lower endpoint is
.056. At probability one it returns a nonempty strictly positive set in
401 of 500 datasets. The corresponding normal-band counts are 193 and
469. Mean invitations rise from about 152 to 1,521, and mean completed
recoveries from about 141 to 1,411. These counts describe information burden,
not bank collection prices or the superiority of either interval rule.

At 94.8\% coverage the Monte Carlo standard error is approximately one
percentage point. The difference from nominal 95\% therefore does not by
itself establish miscalibration; the theorem and its absence for the normal
band remain separate from these empirical proportions. Across all four
invitation rates, raw endpoint crossings occur in 108 intensive-recovery
cells and 712 complete-recovery cells. One complete-recovery normal band is
empty. None is removed from coverage counts, and no empty set is treated
as a supported positive decision. The original signed endpoint separations
remain available alongside the nonnegative set widths used here.


\begin{figure}[htbp]
\centering\includegraphics[width=\textwidth]{partial-recovery-figure.pdf}
\caption{Different purchases of information. Left: exact population effect
sets under the generated completion laws; the dashed line marks the historical
effect of ten percentage points. Right: intensive recovery with different
invitation rates, showing mean confidence-set width over 500 datasets per rate.
The dotted line is the unchanged population-set width. These are generated
comparisons, not bank measurements.}
\end{figure}

\section{A recovered record can concern the wrong outcome}

The fourth procedure completes every invited follow-up but changes the
recorded endpoint for half of the originally unsuccessful, routinely missing
treated customers. It leaves control outcomes unchanged. This deliberately
violates historical-outcome stability. The apparent identified point becomes
.235 although the historical treatment effect remains .10. Inviting more
customers now estimates the wrong contrast more precisely.

At invitation probabilities .10, .25, .50 and 1.00, the bounded sets cover
the historical .10 in 100.0\%, 90.4\%, 0.2\% and 0.0\% of the 500
datasets. The normal bands cover it in 7.0\%, 0.2\%, 0.0\% and 0.0\%.
The wide first band's accidental coverage is not evidence that the stability
assumption is dispensable. At full invitation the bounded construction
covers the different recorded-outcome point .235 in all 500 datasets,
while the normal band covers it in 472. Neither statement validates
inference for the historical target.

A sensitivity analysis must restrict the error that actually matters.
Let $k_a=E\{(1-R)S\mid A=a\}$ and, for $k_a>0$, let
$\widetilde\mu_a$ be the mean recorded outcome among recoverable customers.
A zero recovered mass contributes zero without defining that mean.
Suppose external information
justifies an absolute bound $\gamma_a$ on the difference between that mean
and their historical-outcome mean. The recovered contribution then lies in
\[
 k_a\max(0,\widetilde\mu_a-\gamma_a)
 \ \leq\ E\{(1-R)SY\mid A=a\}\ \leq\
 k_a\min(1,\widetilde\mu_a+\gamma_a).
\]
Adding the routine contribution and unresolved mass gives a sensitivity set
by the same arm subtraction as~\eqref{rel73:identified-set}. This is a
population sensitivity calculation. Estimating its endpoints would require
an additional sampling argument; the original confidence formula cannot be
attached without checking the altered scores and assumptions.

In the unstable construction, $k_1=.38$, the recovered treated success mass
is .245, and its historical value is .11. If control recovery is stable,
the smallest treated mean-drift bound that admits the truth is
$\gamma_1=.135/.38=27/76$, approximately .3553. At $\gamma_1=.30$ the
effect sensitivity set is $[.121,.349]$ and still excludes .10; at .40 it
is $[.083,.370]$. These are stipulated limits on outcome drift, not estimates
learned from the incomplete data.

The distinction matters when choosing software. The inspected
\texttt{attrition} 1.0.0 source at its 21 September 2026 revision implements
the paper's double-sampling bounds, analytic variance and normal critical-value
calculation. Its package description calls the sensitivity analysis a check
on outcome stability, but the function documentation and
\texttt{gen\_mean} implement a different restriction: a fraction $1-\delta$
of follow-up nonresponders has the responder mean, while the remaining
fraction $\delta$ receives worst-case outcomes. That is the paper's
nonresponse-mixture sensitivity, not a bound on historical-outcome drift.
When recovery is complete, there are no follow-up nonresponders and changing
$\delta$ cannot alter the erroneous .235 contrast. Static inspection
establishes this distinction; it is not a native package trial.
\citep{rel73attrition,rel73coppock2017}

\section{The information purchase and its alternatives}

The available implementation is a serious comparator for bounded randomized
experiments with the declared follow-up design. Its sample-ratio estimator,
asymptotic variance and poststratification should be compared under matched
sampling rules and targets. The proposed inverse-invitation calculation makes
known design probabilities and a conservative finite-sample guarantee explicit,
but its crossed endpoints and wide intervals are costs of that choice. Neither
route validates the other merely because their population limits agree.
The package declares GPL-3; a proposed integration must retain its version,
applicable terms and dependency record. No enabled integration or measured
runtime advantage is claimed. \citep{rel73attrition}

Other options change the information or assumptions. A reliable administrative
linkage may recover the fixed historical outcome more credibly than a later
customer report. An intensive protocol can reduce the unresolved mass, while
more invitations to the same protocol mainly improve precision. A model that
assumes completion is ignorable conditional on recorded covariates can support
a point estimate, but it must justify that assumption and diagnose overlap.
The package's trimming option instead targets customers whose outcomes would
be observed under either treatment, with selection assumptions explicitly
chosen. That always-observed target can be useful, but it is not the full
experimental-population effect in this section. Equal observed response rates
alone do not establish individual response monotonicity or justify changing
the target without saying so. \citep{rel73attrition}

To connect identification with a financial decision, stipulate a benefit of
USD~100 for each additional historical repayment and an offer cost of USD~5
per customer. These invented values define net benefit $100\tau-5$, not
observed bank returns. The standard-recovery population set permits net
values from USD~$-13.70$ to USD~10.40. Intensive recovery narrows this to
USD~0.60 through USD~6.08; complete stable recovery gives USD~5.00. Those
are population implications. A finite-sample spending recommendation needs
its own nonempty net-benefit confidence set above zero, and follow-up costs
must also be counted. A positive effect alone is insufficient to clear the
USD~5 threshold.

At full invitation under intensive recovery, only 5 of 500 bounded sets
and 22 normal bands have a nonempty net-benefit range wholly above zero.
The earlier counts of 401 and 469 concerned a positive effect, before
deducting the offer cost. Under complete stable recovery, the corresponding
net-positive counts are 335 and 432. These are decision frequencies in the
same retained datasets. They show how a small population margin can remain
hard to distinguish from zero even after every missing customer is invited. The normal-band frequencies retain their empirical qualification;
more affirmative decisions alone do not establish a better procedure.

The resulting design choice is concrete. The bank should establish whether
follow-up recovers the fixed historical endpoint, record every invitation and
completion, and compare the cost of reducing the unresolved mass with the
cost of estimating a fixed set more precisely. A workflow that cannot defend
stability should expose an appropriately justified sensitivity range or
withhold the historical-effect recommendation. The generated single-treatment
example shows why more invitations, more completed responses and more stable
historical information can have different decision values. Extending that
conclusion to a bank requires evidence about its actual recovery process.
Unbounded longitudinal rewards require their own identification and sampling
arguments, and operating acceptance remains unestablished.

